\documentclass[10pt,a4paper]{amsart}
\usepackage{amsmath,amssymb,amsthm,graphicx}
% Retained preamble for the Stacks Project chapter morphisms.tex (Morphisms of Schemes).
% The chapter file pinned-original-source.tex carries only \input{preamble} on its first line, so no
% preamble text lives inside the chapter itself.  The Stacks Project uses one single shared preamble for
% every chapter; the snapshot stages it verbatim in sources/stacks-chow/preamble.tex.  The block below is
% copied mechanically, line for line, from that file's lines 14-18 (the Xy-pic setup) and lines 154-217
% (the native theorem declarations together with the mathematical macros).  Nothing was added, removed or
% redefined.  The wrapper class and the amsmath/amssymb/amsthm/graphicx base packages are supplied by the
% assembler, so the class/\externaldocument part of the shared preamble is not repeated here.
% For commutative diagrams we use Xy-pic
\usepackage[all]{xy}

% We use 2cell for 2-commutative diagrams.
\xyoption{2cell}
% Theorem environments.
%
\theoremstyle{plain}
\newtheorem{theorem}[subsection]{Theorem}
\newtheorem{proposition}[subsection]{Proposition}
\newtheorem{lemma}[subsection]{Lemma}

\theoremstyle{definition}
\newtheorem{definition}[subsection]{Definition}
\newtheorem{example}[subsection]{Example}
\newtheorem{exercise}[subsection]{Exercise}
\newtheorem{situation}[subsection]{Situation}

\theoremstyle{remark}
\newtheorem{remark}[subsection]{Remark}
\newtheorem{remarks}[subsection]{Remarks}

\numberwithin{equation}{subsection}

% Macros
%
\def\lim{\mathop{\mathrm{lim}}\nolimits}
\def\colim{\mathop{\mathrm{colim}}\nolimits}
\def\Spec{\mathop{\mathrm{Spec}}}
\def\Hom{\mathop{\mathrm{Hom}}\nolimits}
\def\Ext{\mathop{\mathrm{Ext}}\nolimits}
\def\SheafHom{\mathop{\mathcal{H}\!\mathit{om}}\nolimits}
\def\SheafExt{\mathop{\mathcal{E}\!\mathit{xt}}\nolimits}
\def\Sch{\mathit{Sch}}
\def\Mor{\mathop{\mathrm{Mor}}\nolimits}
\def\Ob{\mathop{\mathrm{Ob}}\nolimits}
\def\Sh{\mathop{\mathit{Sh}}\nolimits}
\def\NL{\mathop{N\!L}\nolimits}
\def\CH{\mathop{\mathrm{CH}}\nolimits}
\def\proetale{{pro\text{-}\acute{e}tale}}
\def\etale{{\acute{e}tale}}
\def\QCoh{\mathit{QCoh}}
\def\Ker{\mathop{\mathrm{Ker}}}
\def\Im{\mathop{\mathrm{Im}}}
\def\Coker{\mathop{\mathrm{Coker}}}
\def\Coim{\mathop{\mathrm{Coim}}}

% Boxtimes
%
\DeclareMathSymbol{\boxtimes}{\mathbin}{AMSa}{"02}

%
% Macros for moduli stacks/spaces
%
\def\QCohstack{\mathcal{QC}\!\mathit{oh}}
\def\Cohstack{\mathcal{C}\!\mathit{oh}}
\def\Spacesstack{\mathcal{S}\!\mathit{paces}}
\def\Quotfunctor{\mathrm{Quot}}
\def\Hilbfunctor{\mathrm{Hilb}}
\def\Curvesstack{\mathcal{C}\!\mathit{urves}}
\def\Polarizedstack{\mathcal{P}\!\mathit{olarized}}
\def\Complexesstack{\mathcal{C}\!\mathit{omplexes}}
% \Pic is the operator that assigns to X its picard group, usage \Pic(X)
% \Picardstack_{X/B} denotes the Picard stack of X over B
% \Picardfunctor_{X/B} denotes the Picard functor of X over B
\def\Pic{\mathop{\mathrm{Pic}}\nolimits}
\def\Picardstack{\mathcal{P}\!\mathit{ic}}
\def\Picardfunctor{\mathrm{Pic}}
\def\Deformationcategory{\mathcal{D}\!\mathit{ef}}
\title{A finite birational morphism of integral schemes with normal target is an isomorphism}
\author{The Stacks Project Authors}
\date{stacks-project:morphisms.tex:local-snapshot-20261012; GNU FDL 1.2 or later}
\begin{document}
\maketitle
\noindent\textit{Editorial note.} All delivered mathematical text is the author's own, copied from the pinned Stacks Project chapter file morphisms.tex; no statement or proof step was rewritten, reordered or removed.\par\medskip
\noindent\textit{Editorial note.} The canonical Stacks tag of the counted statement is 0AB1 (morphisms-lemma-finite-birational-over-normal). Following the multi-claim convention for this source, source\_url is the source-level URL https://stacks.math.columbia.edu/ and the tag is carried in claim\_key and here rather than in a per-claim tag URL.\par\medskip
\noindent\textit{Editorial note.} The statement is self-contained: integral scheme, normal scheme, finite and integral morphism and birational morphism are all standard and none of them is introduced by an immediately preceding definition of this chapter that the statement quotes, so no context span is attached.\par\medskip
\noindent\textit{Editorial note.} The parenthetical (or even integral) in the statement is the author's own remark that finiteness is used only through integrality; it is kept verbatim and the proof indeed uses integrality alone.\par\medskip
\noindent\textit{Editorial note.} The preamble is the chapter's shared Stacks Project preamble, retained as sources/stacks-morphisms/retained-preamble.tex: a line-for-line excerpt of sources/stacks-chow/preamble.tex lines 14-18 and 154-217. No macro was invented or redefined.\par\medskip
\noindent\textit{Allowed standard prerequisites (editorial).}\begin{itemize}
\item Integral closure and the definition of a normal scheme. Algebra, Section labelled algebra-section-normal (canonical Stacks tag 00G3); used at proof lines 14640-14641. Supplies that for a normal scheme the ring of sections O\_Y(V) of an affine open is integrally closed in its fraction field, which is the step that forces O\_X(U) = O\_Y(V).
\item A finite morphism is integral, and integrality is checked on an affine open cover. Morphisms, Lemma labelled morphisms-lemma-integral-locally-finite and the definition of an integral morphism in the same chapter; used at proof lines 14641-14643. Gives that every element of O\_X(U) is integral over O\_Y(V), i.e. lies in the integral closure of O\_Y(V).
\item A birational morphism of integral schemes is an isomorphism over a dense open, so the induced map on function fields is an isomorphism. Morphisms, Section labelled morphisms-section-birational (canonical Stacks tag 01RJ); used at proof line 14639. Supplies that O\_Y(V) -> O\_X(U) is injective and induces an isomorphism of fraction fields, so both rings sit inside one common field.
\end{itemize}\medskip
\begin{lemma}
\label{lemma-finite-birational-over-normal}
A finite (or even integral) birational morphism $f : X \to Y$
of integral schemes with $Y$ normal is an isomorphism.
\end{lemma}

\begin{proof}
Let $V \subset Y$ be an affine open
with inverse image $U \subset X$ which is an affine open too.
Since $f$ is a birational morphism of integral schemes, the homomorphism
$\mathcal{O}_Y(V) \to \mathcal{O}_X(U)$ is an injective map of domains
which induces an isomorphism of fraction fields. As $Y$ is normal,
the ring $\mathcal{O}_Y(V)$ is integrally closed in the fraction field.
Since $f$ is finite (or integral) every element of $\mathcal{O}_X(U)$
is integral over $\mathcal{O}_Y(V)$. We conclude that
$\mathcal{O}_Y(V) = \mathcal{O}_X(U)$. This proves that $f$ is an
isomorphism as desired.
\end{proof}

\end{document}
